\section{Propuesta de investigaci\'on}

\subsection{Uso de datos de Aduanet e INEI}

La investigaci\'on utilizar\'a datos de exportaci\'on aduanera (Aduanet) e indicadores agropecuarios del INEI para construir series temporales de precio FOB, volumen exportado, destinos y variables de contexto productivo. La Tabla~\ref{tab:fuentes_datos} resume las fuentes y variables principales. Aduanet aporta granularidad comercial alineada con la realidad exportadora; INEI complementa con producci\'on regional y valor bruto de producci\'on.

\begin{table}[H]
\centering
\caption{Fuentes de datos: Aduanet e INEI}
\label{tab:fuentes_datos}
\small
\begin{tabularx}{\textwidth}{p{2cm}p{3cm}X}
\toprule
\textbf{Fuente} & \textbf{Variables} & \textbf{Descripci\'on} \\
\midrule
Aduanet & Precio FOB, volumen, destino & Exportaciones aduaneras de palta Hass (partida 080440) \\
INEI    & Producci\'on regional, VBP & Estad\'isticas agropecuarias por departamento \\
\bottomrule
\end{tabularx}
\parbox{0.94\textwidth}{\small \textit{Nota.} Elaboraci\'on propia.}
\end{table}


El procesamiento incluir\'a limpieza de registros, tratamiento de valores faltantes, construcci\'on de rezagos, variables estacionales y agregaciones semanales compatibles con horizontes de predicci\'on de 4, 6 y 8 semanas \citep{ReisFilho2022, Torres2022}. La calidad del pipeline de datos condiciona directamente la validez comparativa entre SARIMA, LSTM y SVR.

\subsection{Comparaci\'on entre SARIMA, LSTM y SVR}

Se implementar\'an y comparar\'an tres familias de modelos: SARIMA (l\'inea base estad\'istica), LSTM (aprendizaje profundo) y SVR (aprendizaje autom\'atico con kernels) \citep{Kurnadipare2025, Paul2022, Xia2024}. La evaluaci\'on usar\'a m\'etricas RMSE, MAPE y $R^{2}$ sobre particiones temporales out-of-sample, evitando fuga de informaci\'on futura. La Tabla~\ref{tab:variables_modelo} define variables del modelo predictivo.

\begin{table}[H]
\centering
\caption{Variables del modelo predictivo}
\label{tab:variables_modelo}
\small
\begin{tabularx}{\textwidth}{p{2.5cm}p{2cm}X}
\toprule
\textbf{Variable} & \textbf{Tipo} & \textbf{Descripci\'on} \\
\midrule
Precio FOB & Dependiente & Precio unitario de exportaci\'on (USD/kg) \\
Volumen exportado & Independiente & Toneladas exportadas por periodo \\
Mes / temporada & Independiente & Estacionalidad y ventana comercial \\
Rezago precio & Independiente & Valores rezagados de la serie FOB \\
Destino principal & Independiente & Concentraci\'on por mercado \\
\bottomrule
\end{tabularx}
\parbox{0.94\textwidth}{\small \textit{Nota.} Elaboraci\'on propia.}
\end{table}


La comparaci\'on no busca \'unicamente minimizar error, sino identificar el modelo con mejor equilibrio entre precisi\'on, estabilidad e interpretabilidad para usuarios no t\'ecnicos \citep{Theofilou2025}. Referencias como \citet{Kurnadipare2025} y \citet{Xia2024} sugieren ventajas de LSTM en no linealidades, mientras SARIMA permanece competitivo en series estacionales estables y SVR en conjuntos de tama\~no moderado con variables explicativas.

\subsection{Predicci\'on de precios FOB de 4, 6 y 8 semanas}

Los horizontes de 4, 6 y 8 semanas se alinean con decisiones operativas de campa\~na: programaci\'on de cosecha, negociaci\'on preliminar, ajuste de insumos y evaluaci\'on de viabilidad econ\'omica inmediata. Horizontes m\'as largos aumentan incertidumbre; m\'as cortos limitan utilidad para planificaci\'on de siembra de mediano plazo. La combinaci\'on propuesta cubre necesidades t\'acticas del productor liberte\~no.

Las predicciones alimentar\'an el simulador de escenarios y el m\'odulo de rentabilidad, traduciendo puntos de pron\'ostico en rangos optimista/esperado/pesimista seg\'un bandas de error hist\'orico del modelo seleccionado \citep{Velasquez2025}.

\subsection{Backend en FastAPI}

El backend se desarrollar\'a en FastAPI (Python), exponiendo endpoints REST para entrenamiento/inferencia de modelos, consulta de series hist\'oricas, simulaci\'on de escenarios y generaci\'on de reportes. FastAPI ofrece tipado, documentaci\'on autom\'atica (OpenAPI) y desempe\~no adecuado para servicios anal\'iticos que consumir\'a la aplicaci\'on m\'ovil.

La arquitectura separa capa de datos (ETL Aduanet/INEI), capa de modelos (SARIMA, LSTM, SVR) y capa de negocio (rentabilidad, recomendaciones), integradas mediante una API REST consumida por la aplicaci\'on m\'ovil.

\subsection{Aplicaci\'on m\'ovil en React}

El frontend m\'ovil se implementar\'a en React (React Native o React con empaquetado m\'ovil), consumiendo la API FastAPI para mostrar pron\'osticos, hist\'oricos interactivos, simulaciones y recomendaciones de siembra. La interfaz priorizar\'a legibilidad en condiciones de campo: tipograf\'ia clara, gr\'aficos simples, lenguaje no t\'ecnico y flujos cortos para consultar escenarios.

La Tabla~\ref{tab:product_backlog} resume funciones principales del software seg\'un el Product Backlog del proyecto. La evaluaci\'on con productores de La Libertad (Cap\'itulo II) medir\'a utilidad percibida y mejora en comprensi\'on de rentabilidad.

\begin{table}[H]
\centering
\caption{Product Backlog: funciones principales del software}
\label{tab:product_backlog}
\small
\begin{tabularx}{\textwidth}{p{1.2cm}p{3.5cm}X}
\toprule
\textbf{ID} & \textbf{Funci\'on} & \textbf{Descripci\'on} \\
\midrule
F-01 & Pron\'ostico FOB & Predicci\'on de precios a 4, 6 y 8 semanas \\
F-02 & Simulador & Escenarios optimista, esperado y pesimista \\
F-03 & Rentabilidad & C\'alculo de ingresos, costos y margen \\
F-04 & Hist\'orico & Visualizaci\'on de series y tendencias \\
F-05 & Recomendaci\'on & Sugerencia de siembra seg\'un escenario \\
\bottomrule
\end{tabularx}
\parbox{0.94\textwidth}{\small \textit{Nota.} Elaboraci\'on propia.}
\end{table}


\subsection{Simulador de rentabilidad y recomendaciones de siembra}

El simulador integra predicciones FOB, costos editables por el productor, rendimiento esperado y par\'ammetros log\'isticos para calcular ingresos, costos totales y margen proyectado \citep{Capunay2025}. Las recomendaciones de siembra derivan de umbrales configurables: por ejemplo, si el margen esperado en escenario pesimista es negativo, sugerir postergar expansi\'on de hect\'areas.

La Tabla~\ref{tab:objetivos_solucion} relaciona soluciones t\'ecnicas con tipos anal\'iticos (descriptiva, predictiva, prescriptiva). Las soluciones S-1 a S-5 articulan un sistema integral que va m\'as all\'a del pron\'ostico aislado, respondiendo a limitaciones identificadas en la literatura y en la realidad productiva de La Libertad \citep{VargasCordova2025, Theofilou2025}.

\begin{table}[H]
\centering
\caption{Objetivos de la soluci\'on: Soluci\'on T\'ecnica y Tipo Anal\'itica}
\label{tab:objetivos_solucion}
\small
\begin{tabularx}{\textwidth}{p{2cm}p{3.5cm}X}
\toprule
\textbf{ID} & \textbf{Tipo anal\'itica} & \textbf{Soluci\'on t\'ecnica} \\
\midrule
S-1 & Descriptiva / predictiva & M\'odulo de c\'alculo de rentabilidad \\
S-2 & Prescriptiva & Simulador de escenarios de siembra \\
S-3 & Predictiva & Modelo SARIMA, LSTM y SVR para precios FOB \\
S-4 & Descriptiva & An\'alisis hist\'orico interactivo \\
S-5 & Descriptiva & Reportes de ciclos productivos \\
\bottomrule
\end{tabularx}
\parbox{0.94\textwidth}{\small \textit{Nota.} Elaboraci\'on propia.}
\end{table}

